\documentclass[12pt]{scrreprt}
\usepackage{tocloft}
\usepackage{xcolor}
\usepackage{hyperref}
\usepackage{lineno}
\usepackage{hyperref}
\usepackage{graphicx}

\definecolor{stevensred}{rgb}{0.64, 0.15, 0.22}

\setcounter{tocdepth}{1} %Adjust to change what level headings are shown in table of contents
\setcounter{secnumdepth}{5}

%Setup display of section headings in article
\renewcommand\thechapter{Article \Roman{chapter}}
\renewcommand\thesection{Section \arabic{chapter}.\arabic{section}}
\renewcommand\thesubsection{\Alph{subsection}}
\renewcommand\thesubsubsection{\alph{subsubsection}}
\renewcommand\theparagraph{\arabic{paragraph}}
\renewcommand\thesubparagraph{\alph{subparagraph}}

%Setup display of section headings in referrences and table-of-contents
\makeatletter
\renewcommand{\p@chapter}{Article \Roman{chapter}\expandafter\@gobble}
\renewcommand{\p@section}{Section \arabic{chapter}.\arabic{section}\expandafter\@gobble}
\renewcommand{\p@subsection}{Section \arabic{chapter}.\arabic{section}(\Alph{subsection})\expandafter\@gobble}
\renewcommand{\p@subsubsection}{Section \arabic{chapter}.\arabic{section}(\Alph{subsection})-\alph{subsubsection}\expandafter\@gobble}
\renewcommand{\p@paragraph}{Section \arabic{chapter}.\arabic{section}(\Alph{subsection})-\alph{subsubsection}.\arabic{paragraph}\expandafter\@gobble}
\renewcommand{\p@subparagraph}{Section \arabic{chapter}.\arabic{section}(\Alph{subsection})-\alph{subsubsection}.\arabic{paragraph}.\alph{subparagraph}\expandafter\@gobble}
\makeatother

%Setup padding for headings in table-of-contents
\setlength{\cftchapnumwidth}{6em}
\setlength{\cftsecnumwidth}{6em}
\setlength{\cftsubsecnumwidth}{1.8em}
\setlength{\cftsubsubsecnumwidth}{1.8em}
\setlength{\cftparanumwidth}{1.8em}
\setlength{\cftsubparanumwidth}{1.8em}

%Setup coloring for referrences and table-of-contents
\hypersetup{
  colorlinks=true,
  citecolor=stevensred,
  filecolor=stevensred,
  linkcolor=stevensred,
  urlcolor=stevensred,
  pdfborder={0 0 0}
}

\begin{document}
\title{The Bylaws of the \\ Student Government Association of \\ Stevens Institute of Technology}
\subtitle{
    \vspace{1cm}
    \includegraphics[width=0.6\textwidth]{logo}
}
\date{Effective 2019-03-??}
\maketitle

%auto-generate table-of-contents
\tableofcontents

%start displaying line numbers
\modulolinenumbers[1]
\linenumbers

\chapter{Formation of Law}

\section{Legislation}

\subsection{Definition}
Legislation is a written proposal brought before the Senate by a Senator to take certain action or adopt, remove, or amend certain policies.

\subsection{Types}
Senators can introduce only two types of legislation to the Senate: resolutions and proclamations. A resolution can adjust fund allocations, modify governing documents of the SGA, or it can mandate action of SGA Officers. A proclamation is a decree, declaring an opinion of the Senate; it may address a specific party. Any piece of legislation that is under the consideration of the Senate shall be referred to as a bill. Any passed piece of legislation shall be referred to by its respective type: resolution or proclamation.

\section{Origin of a Bill}
A senator shall submit a bill to the Secretary, who shall place the bill on the docket. A motion to approve the bill must appear on the docket for the Senate to consider the bill. The Senate can approve only bills on the docket.

\section{Structure} \label{sec:bill_content}
\subsection{General}
A piece of legislation shall contain:
\begin{itemize}
    \item a title;
    \item the proposal date
    \item the proposer(s);
    \item the sponsor(s), if any;
    \item the reasons for the legislation, in the form of a “whereas” clauses; and
    \item the intended actions, in the form of “resolved” clauses.
\end{itemize}
    

\subsection{Title}
The proposer(s) of legislation shall title their legislation in accordance to the following
structure:
\begin{enumerate}
    \item The type of legislation, either “Proclamation” or “Resolution”;
    \item An index code with a hyphen (“-”) separating each of the following subitems:
    \begin{enumerate}
        \item “S” for Resolution and “P” for Proclamation;
        \item The last two digits of the current year — i.e. “17” for 2017 — followed by the abbreviation for the semester in which the proposer(s) presents the legislation — i.e. “S” for Spring or “F” for Fall;
        \item A three-digit ordinal number, incrementing for each bill presented to the Secretary; this number shall be prefixed with zeros to the meet three-digit requirement when necessary — i.e. “005” for the fifth bill introduced into the Senate;
    \end{enumerate}
    \item The phrase “of the”, followed by the current Senate session number, followed by “Senate of the Student Government Association of Stevens Institute of Technology”;
    \item The name of the legislation, as determined by the proposer(s).
\end{enumerate}

\section{Proposition}
A bill must have at least one proposer. Only senators shall be proposers of a bill. The proposer(s) of a bill should only be those senators who were involved in the drafting, creating, or collaboration of the bill. A Senator may withdraw themselves as a proposer of the bill at any time until voting of approval begins on the bill. A Senator shall not be added to the list of proposers of the bill once legislation has been added to the docket. Should all proposers withdraw support of a bill The Senate may reserve the right with 2/3 majority to approve the bill with “Senate” as the proposer.

\section{Sponsorship}
A bill may have a list of sponsors. Any relevant person or entity may be listed as a sponsor of a bill. A sponsor of a bill should only be a person or entity who supports the bill, who has an intimate involvement with the bill, and who consents to sponsoring the bill. A person or entity may withdraw their sponsorship of the bill at any time until voting of approval begins on the bill. A person or entity shall not be added to the list of sponsors of the bill once the bill has been added to the docket.


\chapter{The Senate}

\section{Senators}
\subsection{Definition}
Any representative elected by the Student Body to the Senate shall be known as a Senator.
\subsection{Senator Requirements}
A Senator will be responsible for meeting the following requirements:
\begin{enumerate}
    \item A Senator must attend all Meetings of the Senate in a semester, provided, however, that they may accrue up to two absences.
    \begin{enumerate}
        \item Any Senator who partially attends a meaning, defined as arriving after the beginning of business or leaving prior to the completion of business, shall accrue half an absence.
        \item Any Senator who intends to be absent from a meeting shall provide written notice to the Secretary by 24 hours prior to the beginning of the meeting, and the Senator shall only accrue half an absence.
        \item Any Senator who is absent from a meeting without having provided written notice to the Secretary by 24 hours prior to the beginning of the meeting shall accrue one whole absence.
    \end{enumerate}
    \item A Senator must sit on at least one Standing Committee and meet the member requirements of that committee.
    \item A Senator must attend at least five SGA-funded events per semester, provided that those events fall in at least three different months during that semester. Senators must also
attend at least one Oversight Committee approved administrative outreach event per semester. The Senator must report about each event via the Accountability Form within 72 hours of the event start time.
    \item A Senator must send an email to their district at least twice per month.
    \item A Senator must send an email to their district within 72 hours when mandated by the Public Relations Committee members and chair.
    \item A Senator must work on initiatives and pass a monthly Status Update.
    \begin{enumerate}
        \item An initiative shall be defined as a goal that benefits the student body.
        \item The “Status Update” shall consist of:
        \begin{enumerate}
            \item Evidence that the Senator has completed or attempted to complete their initiative of the current month.
            \item A declaration of a new initiative or redeclaration of a current initiative for the next month.
        \end{enumerate}
        \item A Senator must submit a Status Update to the Oversight Committee for approval the last seven days before the end of the month.
        \begin{enumerate}
            \item The Status Update shall require a majority affirmative vote by the members of the Oversight Committee to pass.
            \item If the Status Update fails to pass, then the Senator may re-submit their Status Update within three days of notification of their failure to pass.
            \item If a Senator fails to have their Status Update passed by the Oversight Committee after the fourth day of the month, the failure is final.
        \end{enumerate}
    \end{enumerate}
    \item A Senator must attend the “Leadership Conferences”, or an equivalent event as defined by
the Oversight Committee, once per semester unless explicitly excused by the Oversight
Committee. The Oversight Committee must announce the date and time of the event a
minimum of one month prior to its occurrence, otherwise, senators will no longer be required to
attend the event, though are encouraged to attend if possible.
\end{enumerate}

\section{Senate Districts}
The members of each class year shall be divided into a number of groups equal to the number of Senate seats for that class. These groups shall be equal or approximately equal in size and have randomly determined constituents. Each Senator shall be assigned one of these groups, hereafter referred to as “Districts”, at the beginning of each academic year. Freshmen Senators shall be assigned Districts upon taking office. If a vacancy of a Senate seat is to occur, the Senator who fills the vacancy shall be assigned the same District as the person who previously held the seat. Senators shall be responsible for regular communication with the constituents of their District and shall serve as a means of the constituents of their District sharing their concerns and opinions with the SGA as a whole. The Secretary shall be responsible for organizing and assigning Districts. The term “class year” shall be defined by the Constitution of the SGA.

\section{Elections of the Senate}

\subsection{Frequency}
Elections of the Senate shall occur once per academic year. 

\subsection{Nominations}
Nominations shall open two weeks before the constitutionally mandated election 
period and shall close one week thereafter. A prospective candidate must 
deliver a petition with signatures from 50 students within their class year to 
the Secretary in order to become nominated.

\subsection{Verification} \label{sec:election_verification}
The Secretary, in conjunction with the Office of Student Life, shall then have 
one week to determine if the nominees meet the eligibility requirements to be 
a Senator as defined by Institute policy or established by resolution of the 
Senate. 

\subsection{Election}
The appropriate administrators shall be directed to enable the voting system 
immediately after the completion of the verification period. The election 
shall remain open for one week. The option to reopen nominations shall always be included in the election poll. In the event that any nominees garner less votes than the option to re-open nominations, those nominees shall not be elected in that election cycle. Of the remaining nominees, those receiving the most votes will fill the vacant seats. Any nominee who failed to be elected shall still have the ability to run in future elections. The elected Senators shall assume office at the beginning of the next meeting of the Senate.

\subsection{Vacancy}
In the circumstance of a vacancy of a Senate seat or receipt by the Secretary of a
Senator’s resignation, a vacancy election shall be held. Once the Election
Committee recognizes the need for a vacancy election, the nomination period must
begin within five days. This period after the acknowledgment of the resignation or
vacancy and before the nomination period begins shall be used by the Public
Relations Committee to prepare a publicity campaign, by the Secretary to contact
the appropriate administrators to set up a vacancy election, and by the Election
Committee to determine the election timeline and announce it to the Senate who
shall announce it to the Student Body. The nomination period will last for an
amount of time determined by the Election Committee of no less than three days
and no more than four days. There shall be a verification period immediately following the nomination period which shall last for an amount of time determined
by the Election Committee of no less than one day and no more than two days. The
appropriate administrators shall be instructed to open the voting period, which shall
last for three days, immediately following the closure of the verification period.
The results must be certified by the verification procedure outlined in Section
7.3(E). Elected Senators shall assume office once verified by the Election
Committee. If the preliminary vacancy election procedure ends with Senate seats left vacant,
the vacancy election procedure shall recommence. These subsequent vacancy
elections periods shall begin no more than three days following the end of the the
previous vacancy election period. There shall be no more than three repetitions of
the vacancy election procedure. Members of the Student Body shall have one vote
per vacancy, but may not vote for the same candidate for more than one vacant
seat.

If a vacancy is to occur while the Senate is not in session, the vacancy election
procedure shall commence once the Senate is in session.

\section{Senate Resignations}
In order to resign, a Senator must submit a formal resignation to the Secretary and
Speaker in written notice. By submitting a formal resignation, a Senator immediately forfeits their Senate seat and shall not be able to resume their position
unless elected to the Senate by the Student Body in a general or vacancy election. Once submitted, a formal resignation may not be withdrawn.
\section{Meetings of the Senate}

\subsection{Regular Meeting}
While the Senate is in full session, the Senate shall hold one meeting per week. These meetings shall be held every Sunday at 7:00 PM local time in room 122 of the Babbio Center, provided, however, that the Speaker, Secretary, or two-thirds of Senators may, in written notice,call a meeting to be held in any other suitable place and at any other suitable time at the Institute that is accessible to the public.
 
\subsection{Emergency Meetings}
The President may call an emergency meeting of the Senate by written notice to the Speaker and the Secretary. In the written notice, the President shall designate the purpose of the emergency meeting, and the intended time and place; no other business shall be considered. The emergency meeting must be held in any suitable place at the Institute that is accessible to the public.

An emergency meeting may also be called by two-thirds of all Senators by written notice to the Speaker, the Secretary, and the President. In the written notice, the Senators shall designate the purpose and the intended time and place of the emergency meeting; no other business shall be considered.

For emergency meetings, the Secretary must prepare agenda and deliver the agenda to all Officials of the SGA at least two hours before the meeting start time; the submission deadline
for emergency meetings shall be three hours before the meeting start time. A majority of all Senators, for emergency meetings, shall constitute quorum.

\subsection{Agenda}
The Secretary shall prepare an agenda (or “docket”) for each regular meeting, showing the time
and place and the business items to be discussed, provided that it adheres to the order of
business set forth in these bylaws. The Secretary shall provide a copy of the agenda to all
Officials of the SGA at least 72 hours before the meeting. For items to be placed on the agenda,
an Official of the SGA must submit their bills, communications, memos, discussion items, or any
other matters to the Secretary by written notice before the submission deadline. All submissions
must be included in the agenda.

\subsection{Submission Deadline}
The submission deadline shall be 11:59 PM local time on Wednesday preceding the regular meeting.
However, if the Secretary receives a submission after the deadline has passed, the Secretary
shall have the power to decide whether or not the submission will be added to the agenda. If the
Secretary adds anything to the agenda after already providing a copy of the agenda to each
Official of the SGA, the Secretary must provide an updated agenda for each Official of the SGA,
notifying them that it has been updated. The Secretary may not add any new items to the
agenda if the start of the meeting will occur within 24 hours.

\subsection{Quorum}
Two-thirds of all Senators shall constitute quorum. Quorum must exist for motions to be
adopted. If, at any time during a meeting, a Senator raises a question as to the presence of quorum, the
Speaker shall forthwith direct the Secretary to count the number of Senators present and shall
announce the result, and these proceedings shall be without debate.
\subsection{Speaker}
The chair of all meetings of the Senate shall be the Speaker, unless otherwise described by
these bylaws. The Speaker shall be elected to a term coinciding with the current session of the
Senate. The Speaker shall be elected during the Senate’s first meeting of a new Session. The
first order of business during this meeting shall be to elect a Speaker. Until the Speaker is
elected, the Vice President of Operations shall be acting Speaker at this meeting. Once a
motion to elect a Speaker has been entertained, the Vice-President of Operations shall accept
nominations from the floor. Only Senators may nominate a Speaker, and each nomination must
be seconded by another Senator. The winning nominee shall, in a vote by secret ballot,
accumulate a plurality of the votes of the Senators present.

If the Speaker resigns, the Senate must hold a vacancy election for a Speaker at the next
general meeting. Until a new Speaker is elected, the Speaker pro-tempore shall be acting
Speaker of meetings. For the vacancy election, only Senators may nominate a Speaker, and
each nomination must be seconded by another Senator. The winning nominee of the vacancy
election shall, in a vote by secret ballot, accumulate a plurality of the votes of the Senators
present. Upon completion of the vacancy election, the Speaker pro-tempore shall no longer hold
the Speaker pro-tempore position, and the new Speaker must appoint a new Speaker
pro-tempore.

\subsection{Speaker Pro-tempore}
If the Speaker cannot preside over a Senate meeting, the duty shall fall to the Speaker
pro-tempore only until the Speaker can preside over the Senate. The Speaker pro-tempore shall
be appointed by the Speaker at the first meeting of a new Session of the Senate, and the
appointment does not have to be approved by the Senate. The term for this position coincides
with a Session of the Senate. The Speaker pro-tempore must be a Senator. 

If the Speaker
pro-tempore resigns from this position before the end of the term, then the Speaker must
appoint a new Speaker pro-tempore at the next meeting.
If the Speaker and Speaker pro-tempore cannot preside over the Senate meeting, the duty shall
fall to the Vice-President of Operations.

\subsection{Secretary Pro-tempore}
If the Secretary is unable to attend all or part of a Senate meeting, the Secretary pro-tempore
from the Senate will fill in to take the minutes for the meeting, or until the Secretary returns. The
Secretary pro-tempore shall be appointed by the Secretary, and the appointment does not have
to be approved by the Senate. The term for this position coincides with a Session of the Senate.
The Secretary pro-tempore must be a Senator. If the Secretary pro-tempore resigns from this
position before the end of the term, then the Secretary must appoint a new Secretary
pro-tempore at the next meeting.

If the Secretary and Secretary pro-tempore are unable to attend all or part of the meeting, the
duty shall fall to the Vice-President of Operations.

\subsection{Order of Business}
All business of the Senate shall be taken up for consideration and disposed in the following
order:
\begin{enumerate}
    \item Call to Order
    \item Public Forum
    \item Approval of Minutes
    \item Report from the Speaker
    \item Reports on Standing Committee
    \item Reports on Ad Hoc Committee
    \item Report on Entertainment Committee
    \item Reports on Committee on Student Interests Subcommittees
    \item Reports from Executive Branch Officials
    \item Read Bills
    \item Unread Bills
    \item Other business submitted to the Secretary
    \item New Business
    \item Adjournment 
\end{enumerate}
Senators shall have the authority to deviate from the order of business as necessary for the
efficient conduct of the meeting. However, the Speaker, at their discretion, may refuse to
entertain motions that deviate from the order of business, provided that the Senate may override
the Speaker’s decision by a majority vote.

\subsection{Reading Period}
Every bill shall be subjected to one reading before it may be approved. The approval of a bill
shall not occur in the same meeting as the first reading.
\begin{enumerate} 
    \item First Reading 
\end{enumerate}
The first reading is when a bill is introduced to the Senate. The first reading shall occur by adopting a motion to read a bill. The motion to read a bill shall be stated as “I move to read
[“bill”].” The motion to read a bill shall have the following characteristics:
\begin{enumerate}
    \item It requires a second.
    \item It shall not be amendable.
    \item It shall be debatable. Debate shall last no longer than five minutes, provided, however,
that debate may end sooner at the discretion of two-thirds of the Senate. Debate may
not be extended.
    \item Adoption of this motion shall not require a vote; instead, adoption of this motion shall
occur upon closure of debate.
    \item It shall not be reconsidered.
\end{enumerate}
Upon adoption of the motion to read a bill, (a) the bill is considered read, and (b) the Secretary
shall place a motion to approve the bill on the agenda for the following meeting of the Senate.

\subsection{Obtaining the Floor and Debate}
\begin{enumerate}
    \item Rights of the Speaker

If the Speaker is a Senator, the Speaker may move or second even while presiding over the
meeting. If the Speaker is not a Senator, the Speaker may not move, second, nor vote.
    \item Rights of the Non-Senators
    
Non-Senators may participate in debate, only when called upon by the Speaker. The Speaker
must prioritize debate of Senators over non-Senators, provided, however, the Speaker may
choose to call on non-Senators over Senators as necessary for the efficient conduct of the
meeting. Non-Senators may not move, second, nor vote.
    \item Privilege of Opening Debate
    
When a Senator makes a motion, that Senator shall, at their discretion, have the privilege of
opening debate. The Speaker must call on that Senator once debate begins on a debatable
motion.

However, when a Senator moves to approve a bill that they did not propose, that Senator shall
not open debate; instead, a proposer of the bill shall, at their discretion, have the privilege of
opening debate. The Speaker must call on a proposer of the bill once debate begins on the
motion to approve the bill.
    \item The Five Minute Rule
    
The Speaker shall have the power to close debate for any debatable motion and bring the
Senate to a vote for the immediate question, provided, however, that the Senate shall be
allowed a minimum of five minutes to debate for any debatable motion. The Senate may appeal
the Speaker’s decision by majority vote.


The Senate may choose to end debate before that five minute period is over by adopting a
motion to previous question. However, the motion to previous question shall not be adopted if
both (a) five minutes have not passed since debate began, and (b) at least one person votes
against the adoption of the motion to previous question.
    \item Informality 

Any person, having obtained the floor, does not (a) have to begin speaking by addressing the
Speaker, nor (b) have to avoid personal references. However, any person, having obtained the
floor, must avoid indecorous language.
    \item Maximum Time for Each Speech
    
Any person, having obtained the floor while a debatable motion is immediately pending, shall
speak no longer than ten minutes unless obtaining permission from the Senate. Such
permission can be given by unanimous consent, or by means of a motion to extend the limits of
debate, which requires a two-thirds vote without debate.
    \item Number of Speeches on the Same Question Per Member Per Day

Any person may speak unlimited times for the same question on the same day — except that in
the case of a motion to appeal, only the Speaker can speak twice (the second time at the close
of debate).
\end{enumerate}
\subsection{One Third Abstention}
Should more than one-third of the senators present abstain on a vote of a debatable motion, the
Senate shall not adopt that motion.
\subsection{Minutes}
The Secretary shall record minutes of all regular and emergency meetings of the Senate
showing the time and place; Officials of the SGA present and any non-Official of the SGA
present, provided that they speak at least once during the meeting; Officials of the SGA absent
and excused; Officials of the SGA absent and unexcused; Officials of the SGA either early or
late; all motions made; the results of every vote including the number of yeas, nays, and
abstains, provided that the vote was not taken by voice; the vote of each member for every
motion, provided that the vote was not taken by voice and was non-secret.

The Secretary must also include transcriptions of all discussion at every meeting of the Senate
in the minutes; however, it does not have to be a verbatim record.

Lack of minutes for a particular meeting or part of a meeting will void any action not recorded,
provided that there was no bad faith involved.

\subsection{Censure}
The Senate may vote to censure a Senator if the Senate deems that the Senator is speaking irresponsibly. A vote of seventy-five percent of the Senate is required. 

Should a Senator be censured, they will surrender their right to speak until the business item has been resolved. They shall not lose their right to vote on the business item. 

\section{Sessions of the Senate} \label{sec:senate_sessions}

\subsection{Definition}
The session of a particular Senate shall begin at the start of the Senate 
meeting following the conclusion of Senate elections, and end at the start of 
the next session of the Senate. A session of the Senate is separate from the
parliamentary session, which shall remain defined as the default as adopted
from the Senate’s parliamentary authority.

\subsection{Senate Number}
Each session of the Senate shall be numbered, beginning with the first session 
of the Student Council in April 1913. The Senate of the 2012-2013 school year 
shall be the ninety-ninth Senate. 

\subsection{Recess}
The Senate may declare, by two-thirds resolution, a period in which it shall 
be considered ``in recess.'' No business may be conducted nor meetings held by 
the Senate while it is in recess, except that a meeting may be held to 
consider a resolution to come back from recess ahead of schedule. 

\subsection{Reduced Session}
During any period in which either class is not held or less than half of the 
Student Body is registered for classes, and the Senate is not in recess, the 
Senate shall be considered to be ``in reduced session.'' Attendance to any 
Senate meetings shall not be mandatory during reduced session. 

\subsection{Full Session}
The Senate shall be considered ``in full session,'' or simply ``in session,'' 
at any time that it is neither in recess nor in reduced session. 

\section{Standing Motions}
Standing motions of the Senate may be made by resolution of the majority of 
the Senate and approval of the President. No standing motion shall be 
inconsistent with or supersede the bylaws or Constitution of the SGA. Standing 
motions shall only have effect until the close of the current session of the 
Senate. No standing motion shall be construed to grant additional powers to 
recognized organizations or the SGA. 

\section{Financial Powers}
The Senate shall have the power to transfer money between SGA accounts and 
Campus Organization accounts. The Senate may not transfer money out of the 
Discretionary Fund held by the Cabinet. The Senate shall have the power to 
spend money out of the main SGA account. However, any motion to spend money must be approved by a simple majority of the Senate where quorum prevails. 

\chapter{The Cabinet}

\section{Presidential Elections}

\subsection{Frequency}
The Presidential Election shall occur once per academic year, during the 
constitutionally mandated election period. 

\subsection{Nominations}
Nominations shall open three weeks before the constitutionally mandated 
election period and shall close one week thereafter. A prospective candidate 
must submit a petition with signatures from five percent of the Student Body 
in order to become nominated. Nominations shall be submitted to the Secretary, 
who shall immediately forward them to the election committee. 

\subsection{Verification}
The election committee, in conjunction with the Office of Student Life, shall 
then have one week to determine if the nominees meet the eligibility 
requirements to be a member of the Cabinet established by resolution or 
standing motion of the Senate or by Student Life policy. 

\subsection{Election}
The appropriate administrators shall be directed to enable the voting system 
immediately after the completion of verification. The election shall remain 
open for one week. The President \& Vice President of Operations Elect shall be announced at 
the beginning of the next meeting of the Senate. They shall assume office and 
be sworn in at noon on the first day of the Spring Semester, in accordance 
with the Constitution of the SGA. 

\section{Archiving}
The Secretary will oversee the archiving of all writs, resolutions, bills, 
proclamations, dockets, addendums, and bylaws in the SGA archives. The Secretary will also be 
responsible for archiving the constitutions, and bylaws of all RSOs, as 
outlined in \ref{sec:document_archival}, and any other material presented to 
the Cabinet by a duly authorized officer of an organization. These archives 
will be available to the Student Body electronically.

\section{Meetings of the Cabinet}
The Cabinet shall meet during the constitutionally mandated period. The 
Cabinet may adopt its own rules for how Cabinet meetings shall run.

\section{Senate Meeting Agendas}
The Secretary will conduct the formation of each and every Senate agenda, 
referred to as ``Docket''. The format will consist of: 
\begin{enumerate}
    \item In order, the following: ``Docket for the,'' the number of the 
          current session of the Senate, ``Senate of the Student Government 
          Association,'' and the date of the Senate meeting at which the 
          agenda is addressed.
    \item A discussion period referred to as “Public Forum” in which members of the Student Body who are not Officials of the SGA in attendance may discuss items they wish to bring to the attention of the Senate.
    \item Reports of the Speaker, President, and each member of the Cabinet.
    \item Committee reports of each committee of the SGA, each sub-committee 
          of the Committee on Student Interests, and the reports of any other 
          sub-committees compelled by standing motion of the Senate to be 
          produced at a Senate meeting. 
    \item Items currently lain on the table by the current session of the 
          Senate. 
    \item The order of business to be conducted, which shall begin with Old 
          Business and conclude with New Business. Within each of order of 
          business section, the docket will consist of: 
        \begin{enumerate}
            \item Any bills to be presented to the Senate during the meeting, 
                  in the form of a motion to approve the bill with a Senate 
                  sponsor of the bill listed as the individual expected, but 
                  not required, to make the motion. 
            \item Any other resolutions to be presented to the Senate during 
                  the meeting, in the form of a motion to approve the 
                  resolution with a Senate sponsor of the resolution listed as 
                  the individual expected, but not required, to make the 
                  motion. 
            \item Any motions regarding financial items outside of the scope 
                  of a resolution, such as approval of a budget or 
                  constitution. 
            \item Any business regarding impeachments either brought by a 
                  member of the Senate or necessitated by the bylaws. 
            \item An indicator that any other motions will be entertained from 
                  the Senate before progressing on to other sections of the 
                  order of business. 
        \end{enumerate}
    \item Any other items required to be placed on the Docket by these bylaws 
          or by standing motion of the Senate.
    \item A discussion period following the completion of all business during which the Senate, Cabinet, and any other members of the Student Body in attendance may discuss any issues they see fit for discussion at a Senate meeting.
\end{enumerate}

\section{Presidential Discretionary Fund}
There shall be an amount of money of 2000 reserved. This amount of money shall be called the “Presidential Discretionary Fund. The President shall be authorized to spend money from this fund without the involvement of the Senate. At no time shall the fund exceed ten times the Student Activity Fee. 

\chapter{Impeachment}

\section{Impeaching an Official}

The Oversight Committee shall impeach an Official by adopting the motion “I move to impeach
[“Name”].” The motion to impeach shall have the following characteristics: 
\begin{enumerate}
    \item It can be made only by a member of the Oversight Committee at an Oversight Committee
meeting.
    \item It can be made only on an Official of the SGA.
    \begin{enumerate}
        \item If a member of the Oversight Committee moves to impeach another member of the
Oversight Committee, then the latter Oversight Committee member shall not have a vote.
The chair must appoint another Senator of that member’s class year as a pro tem member
to vote on that member’s impeachment.
        \item If a member of the Oversight Committee moves to impeach the Speaker, then:
        \begin{enumerate}
            \item The Parliamentarian shall serve as acting chair of the committee while this
motion to impeachment is debated.
            \item The Parliamentarian shall not break a tie if the vote on the motion to impeach is divided.
        \end{enumerate}
    \end{enumerate}
    \item It shall not require a second.
    \item It shall not be amendable.
    \item It shall be debatable.
    \item  Adoption of this motion shall require a simple majority.
    \item An affirmative vote on the motion to impeach can be reconsidered, but not a negative vote.
\end{enumerate}

Upon completion of the vote, the Oversight Committee impeaches the Official by charging them with
breaching the governing documents of the SGA.


\section{Appeal Impeachment Trial}
An Official of the SGA who has been impeached may appeal the decision of the Oversight Committee by
trial before the Senate.

The trial must occur if any senator moves “that the Senate try the impeachment of [“Name”].” The motion
shall have the following characteristics:
\begin{enumerate}
    \item It shall require a second.
    \item It shall not be amendable.
    \item It shall not be debatable.
    \item Adoption of this motion shall require a simple majority.
    \item The impeachment trial must begin immediately after an affirmative vote on the motion.
    \item The motion can not be reconsidered.
\end{enumerate}

The trial must be chaired by the Vice President of Operations. However, if the impeachment trial pertains
to the President or the Vice President of Operations, then the Vice President of Operations shall not chair;
instead, the Parliamentarian shall serve as chair.

The Impeachment Trial shall be conducted according to the following order:
\begin{enumerate}
    \item Presentation of evidence against the accused Official.
    \item Presentation of evidence in defense of the accused Official.
    \item Discussion until all comments have been exhausted.
    \item The chair shall call a vote. A three-fourths vote is required to find an Official of the SGA
responsible of breaching the governing documents of the SGA. This decision may not be
reconsidered. If the vote fails, the decision of the Oversight Committee is overturned.
\end{enumerate}

\section{Disciplinary Action} \label{sec:disciplinary_action}
 The Oversight Committee may bring disciplinary action upon that Official of the SGA if the committee
impeached an Official of the SGA and that decision was not overturned by the Senate.

The motion to bring a type of disciplinary action upon an Official of the SGA may be stated as “I move to
bring [“Disciplinary Action] upon [“Name”].” The motion to bring upon a type of disciplinary action
shall have the following characteristics:
\begin{enumerate}
    \item It can bring upon types of disciplinary action outlined only in Section 4.4.
    \item It shall be made only by a member of the Oversight Committee at an Oversight Committee
meeting.
    \begin{enumerate}
        \item If a member of the Oversight Committee moves to bring disciplinary action upon another
member of the Oversight Committee, then the latter Oversight Committee member shall
not have a vote.
        \item  If a member of the Oversight Committee moves to bring disciplinary action upon the
Speaker, then:
    \begin{enumerate}
        \item The Parliamentarian shall serve as acting chair of the committee while this motion is debated.
        \item The Parliamentarian, not the Speaker, shall break a tie if the vote on this motion is divided.
    \end{enumerate}
    \end{enumerate}
    \item It shall only be made on an Official of the SGA who has been impeached at a prior meeting of the
Oversight Committee.
    \item It shall not require a second.
    \item It shall be amendable.
    \item  It shall be debatable.
    \item It shall require a majority vote to be adopted;
    \item  It shall not be reconsidered.
\end{enumerate}

Upon adoption of the motion to bring a type of disciplinary action on an Official of the SGA, that
disciplinary action is in effect and must be immediately executed. The acting chair of the Oversight
Committee must, in written notice, immediately notify all Officials of the SGA that disciplinary action
has been brought upon an Official of the SGA.

\section{Types of Disciplinary Action}

\subsection{The Acknowledgment Requirement}
The Acknowledgment Requirement requires an Official of the SGA to provide, in written notice,
acknowledgement of breaching the governing documents of the SGA to all Officials of the SGA.

\subsection{Censorship}
Censorship shall bar an Official of the SGA from speaking during debate of motions. This disciplinary
action requires a duration to be specified, which shall not exceed past the end of the current Senate
session.

\subsection{Removal of a Special Title}
This disciplinary action removes a special title from an Official of the SGA. Special titles include
appointed committee chairs and elected committee membership.

\subsection{Removal from Office(s)}
Removal from Office removes an Official of the SGA from their Official position(s). This disciplinary
action shall be effective immediately.

\subsection{Bar from the SGA}
This disciplinary action shall remove an Official of the SGA from their current position and shall prohibit
that person from holding Office in SGA ever again. This disciplinary action shall be effective
immediately and last forever.

\chapter{Recognized Student Organizations} \label{sec:rsos}

\section{Recognized Student Organization}

\subsection{Definition}
A Recognized Student Organization, herein referred to as an ``RSO'', shall be 
defined as a Campus Organization that has obtained recognition from the SGA.

\subsection{The New Organization Process}
Any member of the Stevens Institute of Technology undergraduate student
body may apply to form a new student organization recognized by both the Office of Undergraduate Student Life and the SGA. The process found in Addendum D to this Constitution must be followed in order to acquire this recognition.

\subsection{Duration}
A Campus Organization shall be considered an RSO during the semester in which
it has successfully completed the procedures outlined in The New Organization Process and received full RSO status from the SGA, and any following semesters, even when the Senate is not in full session.

\section{Inactive Organizations}

\subsection{Definition}
An Inactive Organization shall be defined as any approved RSO that fails to 
register with the Cabinet \& the Office of Student Life.

\subsection{Determination}
If an RSO becomes inactive, the Vice President of Student Interests or the Office of Student Life reserves the right to approach members of the organization, to determine if the 
organization is still active outside of the auspices of the SGA.

\subsection{Inactive}
If the number of voting members required for quorum of the organization cannot 
be located by the Vice President of Student Interests, or a majority of located voting members do not affirm that the organization is currently active, then the Vice President of Student Interests shall publicize a new emergency election and shall open nominations for a new executive board of the organization.

\subsection{Inactive or Unable to Reorganize}
Should the organization affirm it is an active RSO, but a new executive board 
is not formed, or the new executive board does not meet the requirements for 
recognition within one week after they take office, the organization shall 
return all SGA funded items and equipment and become ineligible for 
recognition by the SGA until it reorganizes following the process outlined in the New Organization Process addendum to these bylaws.

\section{Nonpartisan Funding}
No Campus Organization may use SGA-provided funds to promote, except as 
provided for by the rights of a free press, any candidates for any election. 
Campus Organizations may promote candidate awareness by providing unbiased 
information on each candidate, or by promoting every candidate equally.

\section{Public Forum}
The publications and content distribution mechanisms of all organizations 
shall be considered a ``public forum.'' Individuals shall be liable for their 
own statements when conveyed by organizations.

\section{Officers and the Executive Board}

\subsection{Unified Elections}

\paragraph{Officer Terms}
Officers of an RSO will serve a term of one-year or until the next
Unified Elections, and until their successor is elected, or a lesser
term if specified in the RSO’s Constitution. Officers of an RSO will
take office immediately following their election, unless otherwise
specified in the RSO’s Constitution.

\paragraph{Timeline}
Unified Elections shall occur once every academic year. The first full
week of school following the University-defined Spring Recess shall be
the Unified Elections period. The last week of school preceding Spring
Recess shall be the Unified Nominations period. An RSO must accept
officer nominations during the Unified Nominations period, submit all
nominations to the Office of Student Life immediately thereafter, and
then hold a special meeting for elections of all officers during the
Unified Elections period, in which any ballot cast for an individual
whose nomination was not previously approved by the Office of Student
Life will be considered invalid.

\paragraph{Custom Elections}
An RSO may define in their Constitution a custom period for their
Unified Elections, provided there is at least one week in between the
Unified Nominations period and the Unified Elections period, and that
the elections process complies with all other rules specified in these
bylaws, especially above.

\paragraph{Changes}
Any changes in an RSO’s officers outside of the standard Unified
Elections period (including RSOs that conduct Unified Elections during
a custom time period) may only occur after notification of the Vice President of Student Interests and after verification by the Office of Student Life of any individuals intending to become an officer.

\paragraph{Vacancies}
Unless the RSO’s Constitution specified otherwise, when a vacancy in
an RSO officer position occurs, a special election process shall be
held following the same procedures as Unified Elections, except that
the Unified Nominations period will begin the first full week of class
following the first meeting after the vacancy exists, and the Unified
Elections period will begin the second full week of class following
the end of the Unified Nominations period.

\subsection{Executive Board}

\paragraph{Definition}
Every RSO will have an Executive Board, consisting only of officers of
the RSO, and including all officers of the RSO unless otherwise
defined in the RSO’s Constitution. A majority of the RSO’s Executive
Board must consist of elected officers.

\paragraph{Funding}
Only a majority of the Executive Board, or higher if specified by the
RSO by Constitution, bylaw, or standing rule, shall have the power to
authorize spending of the RSO’s SGA-allocated funds, provided such
spending is in compliance with the RSO’s budget as approved by the
SGA.

\paragraph{Appointed Officers}
Any officers of an RSO that are not elected by the RSO’s general body
must be appointed by a majority, or higher if specified by the RSO by
Constitution, bylaw, or standing rule, of only those members of the
Executive Board who themselves were elected by the RSO’s general body.

\section{Financing of Campus Organizations}

\subsection{Fund Allocation}
A fund allocation is any amount of money that the Senate transfers between SGA accounts and Campus Organization accounts. The Campus Organization will be responsible for following all policies outlined in The Financial Policy Addendum to these bylaws in order to receive and maintain fund allocations from the SGA.

\section{Votes of Executive Boards}
Any meeting not open to the entire Student Body shall be considered to be held 
in executive session. The executive board of an RSO may only hold exclusive votes on items that are left to solely its discretion that are outlined in the RSO’s constitution. The voting membership of the general body of an RSO shall have the right to vote on all items not left solely to the discretion of its executive board.

\section{Overruling of Quorum or Voting Membership}
Should an RSO become unable to conduct business due to its inability to achieve quorum or the necessary amount of eligible voting members at any of its meetings as a result of an ineffective definition of quorum or voting membership in its constitution or bylaws, the RSO’s quorum or definition of voting membership may be temporarily overruled by the Vice President of Student Interests. The Vice President of Student Interests shall mandate a more suitable, temporary definition of quorum or voting membership for the RSO that allows for business to be conducted at its meetings. This temporary definition of quorum or voting membership shall only last until the RSO has voted on the business it initially sought to conduct, amended its definition of quorum or voting membership in its constitution or bylaws to be more suitable, and had any changes to its constitution reviewed by the Rules Committee and approved by the Senate. This process shall not occur unless the RSO submits evidence to the Vice President of Student Interests of having held at least one special meeting designated solely for the purpose of conducting business and having notified the entirety of its voting membership of this meeting at least one week prior.

\chapter{Committees}

\section{Types}

\subsection{Standing Committees}
Standing Committees shall be defined as permanent committees of the SGA which shall exist to perform certain duties on an ongoing basis. In order for a committee of the SGA to be classified as a Standing Committee, it must have its name and purpose outlined in Article VII of these Bylaws. Standing Committees may have further rules and responsibilities outlined in other governing documents of the SGA, such as the Constitution and addendums to these bylaws.

\subsection{Ad Hoc Committees}
An Ad Hoc Committee shall be defined as a committee of the SGA which shall exist only temporarily to perform certain duties. No governing documents of the SGA shall refer to any Ad Hoc Committee by name.

\section{Chairs}
Unless otherwise specified in any governing documents of the SGA, each committee shall have one chair, which shall be a Senator.

\section{Members}
Unless otherwise specified in the definition of the body, every committee 
shall allow any member of the Stevens community to be a member of the 
committee.

Every Senator, with the exception of the Speaker if they are a Senator, must serve on at least one Standing Committee. Participation on the Committee on Student Interests, Election Committee shall not satisfy this requirement. 
Senators shall be assigned to Standing Committees by the Vice President of Operations based on each Senator’s interest and the need for Senators on each Committee. Senators must meet attendance requirements of each Standing Committee on which they serve as well as perform all duties assigned to them by the Committee Chair, or be eligible for disciplinary action as defined in \ref{sec:disciplinary_action}.

\section{Positions}
The chair of a committee shall be permitted to create or remove officer 
positions in the committee at will, and appoint members to those positions, 
unless otherwise specified in the definition of the committee.

\section{Sub-Committees}
Standing sub-committees shall exist only within the Committee of Student 
Interests, and are further defined in \ref{sec:csi_subcommittees}.

Ad-hoc sub-committees may exist at the discretion of the parent committee. The 
chairs of ad-hoc sub-committees shall be appointed by the chair of their 
parent committee.

\section{Prohibition on Committee Bypass}
Committees may be charged with various duties otherwise left to the entire 
Senate, either by the bylaws or by administrative assignment. Committees may 
also be charged with the duty of prior review before a motion may be presented 
to the Senate. A motion before the Senate that bypasses these duties of a 
committee shall not be entertained, unless the committee has already been 
presented with the issue and has either not resolved it in a timely manner, or 
has ruled in a manner unacceptable to the person or organization presenting 
the issue.

\chapter{Standing Committees}

\section{The Committee on Student Interests} \label{sec:csi_representative}

\subsection{Purpose}
The Committee on Student Interests (“CSI”) shall be the forum for the exchange of
information between Student Organizations and officials of the SGA.

\subsection{Composition}
The CSI shall be composed of the Vice President of Student Interests (“VPSI”), one
representative of each subcommittee, herein referred to as the “Subcommittee Head.”
Responsibilities of Subcommittee Heads can be found in the SHREK addendum. Other
members may be added or removed at the discretion of the VPSI.

\subsection{Standing Sub-Committees of the CSI} \label{sec:csi_subcommittees}
The CSI shall be broken into Standing Subcommittees as determined by these bylaws. Each
student organization shall contribute at least one representative to the composition of their
subcommittee, as defined by their respective constitutions. Each subcommittee shall be
chaired by its respective Subcommittee Head. Each Subcommittee Head shall be elected by
plurality from among the subcommittee member organizations. The Subcommittee Head is
not required to be an SGA Senator.


The election of a Subcommittee Head shall be only called by one of the following: the
Subcommittee Head; the VPSI; or a two-thirds majority of the constituent student organization
representatives. There are no term limits to the position of Subcommittee Head with the
exception of undergraduate commencement.
Each subcommittee shall be required to hold at least one official meeting per month during
the fall and spring semesters.

\subsubsection{Ethnic Student Council}

\paragraph{Purpose}
The purpose of the Ethnic Student Council shall be to to promote unity and good relations among ethnic organizations, to promote and celebrate cultural diversity
among the student body, to encourage membership in ethnic organizations, and to coordinate efforts between the ethnic organizations and the SGA.

\paragraph{Responsibilities}
The Ethnic Student Council shall consider for acceptance, hold, and
execute such responsibilities as delegated to the Council from the member Ethnic RSOs
and the SGA. It shall also hold all responsibilities as an RSO, as dictated by its constitution.

\subsubsection{Professional Societies Subcommittee}

\paragraph{Purpose}
​The purpose of the Professional Societies Subcommittee shall be to coordinate
efforts of Professional RSOs, to coordinate efforts between the administration and
Professional RSOs, to unite Professional RSOs, to coordinate national efforts, to encourage
membership in Professional RSOs, and to stimulate interest in the goals of Professional
RSOs among individuals who are not majoring in the targeted field.


\paragraph{Responsibilities}
The Professional Societies Subcommittee shall consider for
acceptance, hold, and execute such responsibilities as delegated to the Subcommittee from
the member Professional RSOs and the SGA.

\subsubsection{Media Subcommittee}
The ESC Secretary shall maintain accurate and complete minutes of all
meetings. The Secretary shall prepare all correspondence of the Committee.

\paragraph{Purpose}
​The purpose of the Media Subcommittee shall be to coordinate efforts of the
Media RSOs, to collectively protect the rights afforded by the First Amendment of the
Constitution of the United States of America on campus, to unite the Media RSOs, to
encourage membership in Media RSOs, and to stimulate interest in the products of the
Media RSOs.

\paragraph{Responsibilities}
​The Media Subcommittee shall consider for acceptance, hold, and
execute such responsibilities as delegated to the Subcommittee from the member Media
RSOs and the SGA.

\subsubsection{Special Interests Subcommittee}

\paragraph{Purpose}
The purpose of the Special Interests Subcommittee shall be to coordinate
efforts of the Special Interest RSOs, to unite the Special Interest RSOs, and to encourage
membership in Special Interest RSOs.

\paragraph{Responsibilities}
​The Special Interests Subcommittee shall consider for acceptance,
hold, and execute such responsibilities as delegated to the Subcommittee from the Special
Interest RSOs and the SGA.

\subsubsection{Club Sports Subcommittee}

\paragraph{Purpose}
The purpose of the Club Sports Subcommittee shall be to coordinate efforts of
the Club Sports, to coordinate efforts between the administration and the Club Sports, to
unite the Club Sports, to encourage membership in the Club Sports, and to stimulate
interest in the Club Sports.

\paragraph{Responsibilities}
​The Club Sports Subcommittee shall consider for acceptance, hold,
and execute such responsibilities as delegated to the Subcommittee from the member Club
Sports and the SGA.

\subsubsection{Arts \& Music}

\paragraph{Purpose}
The purpose of the Arts and Music Subcommittee shall be to coordinate efforts
of Arts and Music RSOs, to coordinate efforts between the administration and Arts and
Music RSOs, to unite Arts and Music RSOs, and to stimulate interest in the programming
created by Arts and Music RSOs.

\paragraph{Responsibilities}
​The Arts and Music Subcommittee shall consider for acceptance, hold,
and execute such responsibilities as delegated to the Subcommittee from the member Arts
and Music RSOs and the SGA.

\subsubsection{Recreation Subcommittee}

\paragraph{Purpose}
​The purpose of the Recreation Subcommittee shall be to coordinate efforts of
the Recreational Sporting RSOs, to coordinate efforts between the administration and
Recreational Sporting RSOs, to unite, encourage membership, and to stimulate interest in
Recreational Sporting RSOs.

\paragraph{Responsibilities}
​The Recreation Subcommittee shall consider for acceptance, hold and
execute such responsibilities as delegated to the Subcommittee from its member
Recreational Sporting RSOs and the SGA.

\subsubsection{Electronics and Gaming Subcommittee}

\paragraph{Purpose}
​The purpose of the Electronics and Gaming Subcommittee shall be to
coordinate efforts of the Electronics and Gaming RSOs, to coordinate efforts between the
administration and Electronics and Gaming RSOs, and to unite, encourage membership for,
and simulate interest in Electronics and Gaming RSOs.

\paragraph{Responsibilities}
The Sub-committee shall consider for acceptance, hold, and execute such responsibilities as delegated to the Sub-committee from the Electronics and Gaming RSOs and the SGA.

\subsubsection{Service Subcommittee}

\paragraph{Purpose}
​The purpose of the Service Subcommittee shall be to coordinate efforts
between the administration and the Service RSOs, to encourage membership in Service
RSOs, to advertise opportunities and encourage campus to take action and do something
about national and global issues, and to provide a network for the Service RSOs to promote
educational seminars, events, or community service to benefit their specific organization,
but more importantly, the community.

\paragraph{Responsibilities}
The Committee shall consider for acceptance, hold, and execute such
responsibilities as delegated to the Committee from the member Service
RSOs and the SGA.


\subsubsection{Religious and Faith Based Subcommittee}

\paragraph{Purpose}
​The purpose of the Religious and Faith-Based Subcommittee shall be to
coordinate efforts of the religious and faith based RSOs, to unite these RSOs, and to
celebrate and educate about the religious diversity among the Student Body.

\paragraph{Responsibilities}
The Religious and Faith-Based Subcommittee shall consider for
acceptance, hold, and execute such responsibilities as delegated to the Subcommittee from
the Religious and Faith-Based RSOs and the SGA.

\subsubsection{Non-RSO Subcommittee}

\paragraph{Purpose}
​The purpose of the Non-RSO Subcommittee shall be to coordinate the efforts of
student organizations classified as “Non-RSO”s, to unite them, and to promote, protect, and
celebrate their rights as Non-RSOs.

\paragraph{Responsibilities}
The Non-RSO Subcommittee shall consider for acceptance, hold, and
execute such responsibilities as delegated to the Subcommittee by its member Non-RSO
organizations and the SGA.

\subsection{Events hosted by CSI}
Events hosted by the CSI, which shall be called “Leadership Conferences,” shall be
conducted at least once per semester and must be announced at least two weeks prior to
the event. The purpose of these events shall be to extend updates on behalf of the SGA
and the Office of Undergraduate Student Life to Student Organizations, as well as to offer
leadership development to Student Organizations’ executive boards. The VPSI shall
suspend the funding of any Student Organization that fails to send a representative to
events hosted by the CSI.

\subsection{Meetings of the CSI}
The CSI shall meet monthly to hold general discussion regarding Student Organizations.
The CSI shall also be responsible for its duties outlined in The New Organization Process
addendum to these bylaws.

\section{Rules Committee}

\subsection{Purpose}
The purpose of the Committee is to ensure the Constitution, Bylaws, and addendums of the SGA remain up
to date, as well as serve as a resource for Senators wishing to ensure legislation presented to the Senate is properly written and formatted. The Committee must also be responsible for reviewing the constitutions and bylaws of the various RSOs of Stevens Institute of Technology.

\subsection{Review of RSO Constitutions and Bylaws}
The Committee is charged with reviewing any and all constitutions and bylaws of organizations seeking RSO status before they are presented to the Senate, as well as any amendments to an existing RSO’s constitution or bylaws approved by its voting membership. The Committee must insist that any presented constitutions be in compliance with the various rules for RSOs defined in 
\ref{sec:rsos}, be grammatically correct and coherent, and offer other recommendations to ensure for the longevity of the organization. The Committee must ensure the constitution of an RSO, at minimum, includes the name of the organization, a stated purpose, voting membership requirements, a list of executive board positions and associated responsibilities, a stated election period for officers of the executive board, a definition of quorum, and a chaos clause.

The Committee must, upon request, assist organizations in creating appropriate bylaws for the governance of their club, and must bring to the Senate's attention any organization's bylaws
that the Committee feels is in conflict with the organization's constitution, the Constitution of
the SGA, or these bylaws.


Review of RSO governing documents must follow the following procedure:
\begin{enumerate}
    \item Upon receiving RSO governing documents, any information left out of the original
submission will be asked for by a member of the committee.
    \item The submitted governing documents will be reviewed and commented on by at least
two members of the committee, which will be assigned during the soonest possible
Rules Committee meeting.
    \begin{enumerate}
        \item If there are no issues or concerns, the documents must move on to committee approval.
        \item If there are comments, the documents along with comments will be sent back to the original submitter for revisal.
    \end{enumerate}
    \item Revised documents and documents without raised issues or concerns can be
considered for approval through a 2/3 majority vote of the committee where quorum prevails.
    \begin{enumerate}
        \item If the vote passes, the governing documents will be considered approved by the Rules Committee, and be sent to the Senate for ratification or denial.
        \item If the vote fails, the documents will be sent back to the original submitter with comments for further revisal.
     \end{enumerate}
\end{enumerate}


\subsection{Review of Governing Documents}
The Rules Committee must have a reading
of the Constitution of the SGA during the Fall semester and a reading of the Bylaws and
Addendums to these Bylaws during the Spring semester. Readings must completed by the end
of the first full month of classes.If the Committee deems that revisions be necessary to any governing documents of the SGA,the Committee Chair must present at least one bill to the Senate within that semester which must, at a minimum propose revisions which remove any outdated policies and add any new policies not yet codified in each respective document. These bills must be titled “Recommendations of the Rules Committee for the Upkeep of”, followed by either “the Constitution of the Student Government Association”, “the Bylaws of the Student Government Association”, or “the Addendums to the Bylaws of
the Student Government Association” followed by “for the”, followed by the name of the semester in which the bill is being proposed, followed by “Semester”. 

\subsection{Review of Legislation}
The Committee must, upon the request of a Senator, be charged with reviewing any legislation not yet submitted to the Secretary. The Committee shall ensure that the legislation meets the requirements set forth in Article I, is grammatically correct and coherent, does not conflict with the governing documents of the SGA where not intended, and offer other objective assistance to ensure the wording of the legislation is aligned with its purpose.

\subsection{Document Archival} \label{sec:document_archival}
The Committee shall  assist the Secretary in archiving the constitutions and bylaws of RSOs, and the Constitution, bylaws, and addendums of the Student Government Association. The Rules Committee Chair may make non-substantive
changes in spelling, grammar, or numbering to these Bylaws without amendment, provided
they must notify the Secretary and all voting members of the Rules Committee of these
changes in written notice. The Secretary and voting members of the Rules Committee may
object to these changes through written notice to the Committee Chair, and if at least one
objection is raised before 24 hours, the non-substantive changes must be reversed.

\subsection{Quorum}
Majority of all voting members of the Committee shall constitute quorum. Quorum must be
present for motions to be adopted.

\subsection{Voting Membership}
Any undergraduate student of the Stevens Institute of Technology is a voting member of the
Rules Committee if they have attended four out of the past six regularly-scheduled Rules
Committee meetings, or two out of the last four Rules Committee meetings for Officials of the
SGA. The Committee Chair is always a voting member.

\subsection{Meetings}
All meetings of the Rules Committee require an agenda, which must be sent to the voting
members of the Committee at least 24 hours beforehand. This agenda must contain, when
applicable, organization constitutions to be reviewed, suggestions for motions to approve
revised constitutions, and legislation submitted to the Committee for objective review.
Approval of a revised organization constitution requires a two-thirds majority vote where
quorum prevails. The Committee Chair shall be the Chair for the meeting, and designate a
voting member to take minutes for the meeting.


\section{Election Committee}

\subsection{Purpose}
The Committee shall be responsible for administrating all active elections for 
the Cabinet and Senate. The Election Committee shall be responsible for organizing a debate between the candidates/nominees for the Freshman Senator Elections, as well as for the
Presidential Elections. The debates must fall within 2 weeks’ time of the end of the
nomination period. The election committee shall announce the dates of
nomination, dates of elections, and date of the debate to the senate. The senate
shall announce these election rules to the student body. 

\subsection{Requirements of Committee Members}
The term class year shall be interpreted as defined by the Constitution of the 
SGA. In accordance with the Constitution of the SGA, the Secretary shall chair the Election Committee. The Secretary shall appoint one member of each grade from
the Honor Board to serve on the election committee. The sophomore, junior, and
senior representative of the election committee shall be appointed within the first 2
weeks of the academic year. The freshman representative of the election shall be
appointed within 2 weeks of the freshmen honor board elections.

\subsection{Conflict of Interest}
Should a member of the Student Body run for a Cabinet position, be promised an 
appointed Cabinet position, or be running for any Senate seat, the person 
shall become ineligible for membership on the Committee.

\subsection{Improprieties and Errors}
Should any candidate/nominee fail to abide by the rules and procedures of the
election s outlined in The Election Procedures and Guidelines Addendum, the infraction shall be brought up to the Election Committee. Should the Election Committee deem an infraction valid, they shall then decide whether the infraction is Minor or Severe. The Election Committee shall have the right to issue a repercussion listed within the Infraction Repercussion Table. If a desired repercussion strays from the Infraction Repercussion Table then the Election Committee can make a recommendation and only by 3/4 majority, the senate may issue a repercussion that strays from the Infraction Repercussion Table. Consult the Election Procedures and Guidelines Addendum for the possible repercussions. 

\subsection{Result Certification}
The Committee must certify the results of every election of the Senate or 
Cabinet. Should the Committee refuse to certify the results, two-thirds of the 
Senate shall have an opportunity to certify the results by resolution. Should both the Election Committee and the Senate refuse to certify the results, the President shall nominate a member of the student body from each class to certify
the results. 

\section{Public Relations Committee}

\subsection{Purpose}
The purpose of the Public Relations Committee shall be to facilitate and enhance active two-way
communication between the SGA and the Student Body, to inform the community of 
any SGA actions and initiatives, to organize events that promote the SGA, and 
to recognize outstanding achievements by all members of the community. 

\subsection{District Emails}
The Chair of the Public Relations Committee shall appoint appoint a Senator as the District Distribution List Compiler. The District Distribution List Compiler shall compile a list  containing information that they have deemed worthy of sending to the student body. This shall be referred to as the “District Distribution List.” Other
Officials of the SGA may suggest information to the Compiler. The District Distribution List shall be sent
in written notice to all Senators within 26 hours following the start of each meeting of the Senate.

\subsection{Mandated Emails}
Should the Chair of the Public Relations Committee deem any specific piece of information to be crucial to the knowledge of the Student Body, they may mandate that all Senators communicate that piece of information to their Districts within a certain amount of time that shall be no less than seventy-two (72) hours. This requirement can be overturned by a two-thirds majority vote of the Senate. If the Public Relations Committee Chair mandates that a piece of information be communicated with the Student Body before the next Senate meeting is to be held, they must administer an electronic poll such that the Senate may retain its ability to overturn the requirement.



\section{Campus Life Committee}

\subsection{Purpose}
The purpose of the Committee shall be to ensure the demands of the student 
body are represented in all non-academic facets which directly affect the 
quality of life on campus, to provide the opportunity for student input on any 
major policy and service change, to improve the residence life of the student 
body, to improve the student life of commuter students to ensure the dining 
service on campus is meeting the needs and desires of the student body, to 
provide for student input on the operations of all facilities on campus, and 
to define areas of campus improvement that the student body would see as a 
benefit to student life and to recommend the use of the Capital Improvement 
Fund to address these areas.

\subsection{Student Center Liaison}
A member of the Campus Life Committee shall coordinate with Student Life in 
representing student interests for aspects of a student center delegated to 
the students, and any replacement student center constructed or delegated to 
the students.

\section{The Student Faculty Alliance}

\subsection{Purpose}
To unite selected Stevens’ students, faculty, and academic administrators who have
demonstrated diplomatic ability, uncommon foresight and selflessness. These parties
should share the interests of adding value to the Stevens’ education, improving student-faculty
communication, and introducing a common forum for positive progression both
in and out of the classroom.

\subsection{Incorporation}
The Student Faculty Alliance shall operate as a committee of the Student Government
Association, reporting to and managed by the Vice President of Academic Affairs of the
Student Government Association at Stevens Institute of Technology.

\subsection{Chair}
The elected President of the Student Faculty Alliance shall serve as the chair of the
Student Faculty Alliance and shall be elected by all procedures outlined in Addendum
X: The Operations of the Student Faculty Alliance.

\section{Committee on Academic Rights}

\subsection{Purpose}
The purpose of the Committee on Academic Rights is to assist and aid the current
systems on the campus of Stevens Institute of Technology that promote and uphold the
values and rights of students under the Stevens Honor System.

\subsection{The Honor Committee of the Gear and Triangle Honor Society}
The Committee on Academic Rights will work in tandem with the Honor Committee of
the Gear and Triangle Honor Society to promote the Honor System as a whole on the
campus of Stevens Institute of Technology through events, awareness, and other
methods. 

\subsection{The Stevens Honor Board}
The Committee on Academic Rights will serve to assist the Stevens Honor Board in all
endeavors “to promote honor and integrity throughout the Stevens campus both
academically and socially.”

\subsection{Meetings}
The chairperson of the Committee on Academic Rights will propose a meeting at least
twice per semester with representatives from the Committee on Academic Rights, the
Gear and Triangle Honor Committee, and the Honor Board to discuss progress, establish
goals, and create initiatives to facilitate the Stevens Honor System on campus. 

\section{Academic and Curriculum Advancement Committee}

\subsection{Purpose}
The purpose of the Academic and Curriculum Advancement Committee is to represent
the needs and concerns of the Stevens student, and work with the Faculty Senate to
make certain that University educational endeavors are pursued with the student body
in mind.

\subsection{Establishment}
The Academic and Curriculum Advancement Committee will serve under the supervision
of the Vice President of Academic Affairs of the Student Government Association at
Stevens Institute of Technology. 

\subsection{Chairman and Representative to the Faculty Senate}
The chairman of the Academic and Curriculum Advancement Committee will be
appointed by the Vice President of Academic Affairs. This chairman shall serve as the
Student Government Association representative to the Faculty Senate or shall delegate
one Senator to serve as such in their place.

\section{The Senate Budget Committee}

\subsection{Purpose and Powers}
While the power to transfer money between SGA
accounts and Campus Organization accounts is vested in the Senate, the purpose of the Senate Budget Committee of the SGA shall be to make recommendations to the Senate regarding fund allocations. The Senate Budget Committee shall also have the power to independently approve Reallocation Requests. A Reallocation Request allows a campus organization to use fund
allocations from the SGA on items and services different than the original intent.

\section{The Student Affairs Committee}

\subsection{Purpose}
The purpose of the Student Affairs Committee will be to create a corps of liaisons to the
various offices at Stevens Institute of Technology. The offices required to liaison to will be
determined by the Chair of the Student Affairs Committee dependent upon relations already
maintained with these offices and their activities. These liaisons will be responsible for any
feedback or change on behalf of their office and will be regarded the “experts” related to
matters which concern their office.

\subsection{Chair and Membership}
The Chair of the Student Affairs Committee will be the Vice President of Operations.
Liaisons will be appointed by the Vice President of Operations, and will have the responsibility of beginning and continuing a relationship between their respective Office and that of the SGA. Liaisons shall become well versed in the operating procedures, goals, principles, and other aspects of the office that they interface with such that they may effectively help with any feedback or other communications necessary.

\section{The Health and Wellness Committee}

\subsection{Purpose}
The purpose of the Health and Wellness Committee shall be to represent the needs and
concerns of the Stevens student body in regards to their mental, physical, and sexual
wellbeing. The committee shall work with the campus health center and administration to
accommodate students’ requests for health services. The committee shall also work
collaboratively with the Office of Counseling and Psychological Services, the wellness
center, Campus Recreation, the Title IX Coordinator, and other organizations to sponsor
and run events on campus that address areas of mental, physical, and sexual health as
well as act as a liaison between these offices and the student body. The committee shall also serve as a channel for students who are interested in facilitating discussions and workshops that serve to educate the student body on mental, physical, and sexual health to do so. 

\chapter{Addendums to These Bylaws} \label{sec:addendums}
An addendum is any supplementary policy or set of policies adopted by the Senate that is to exist separate from the Constitution of the SGA and these bylaws. No addendum may supersede or be inconsistent with the Constitution of the SGA or these bylaws. 
Addendums shall be adopted after the procedure for passing a bill is followed. 
Amendments to an addendum shall occur after the procedure for passing a bill is followed. An addendum shall be removed from law after the procedure for passing a bill is followed.

\chapter{Amendments to These Bylaws} \label{sec:amendments}
Amendments to these bylaws shall occur after the procedure for passing a bill 
is followed. 

\chapter{Implementation}
Upon ratification of these bylaws, all standing committees and sub-committees 
will be organized as described in these bylaws. These bylaws shall take effect 
immediately after ratification and all further items of business at that 
Senate meeting shall follow these bylaws. 

\chapter{Suspension of These Bylaws and Addendums} \label{sec:suspension}
The motion to suspend any bylaws will have the same characteristics and rules as the motion to
suspend the rules in Robert’s Rules of Order Newly Revised, except that this motion shall be
debatable, can only suspend these Bylaws and addendums to these Bylaws, and requires a two-thirds
vote of the entire membership of the Senate.Article VIII, Article IX, Article X, Article XI, Article XIII, Article XIV, and
Article XV of these bylaws may never be suspended. The motion to suspend the bylaws cannot be
used to suspend any other governing document. Any suspension of the bylaws shall only last until
the end of the Senate meeting. 

\chapter{Miscellaneous}

\section{Written Notice}
Only a physical paper letter or a message sent from a Stevens-issued email account shall satisfy any provision
requiring written notice.

\section{Capital Improvements and Infrastructure}
At least three percent of the total amount of student activity fees collected 
by the SGA shall be saved for capital expenditures benefiting the campus. 
These expenditures need not occur during the semester the money is collected, 
but the money set aside shall not be eligible for allocation for any purposes 
other than capital improvements during later semesters. 

However, if the Capital Improvements Fund has accumulated more than fifty 
percent of the Student Activities Fund from the most recent fall semester, 
only two percent of the collected student activity fees shall be diverted to 
the Capital Improvements Fund, unless the Senate passes a standing motion to 
the contrary. 

\section{Headings}
The heading under which text falls, such as section title or article name, 
shall not bound the effect of that content. 

\section{Senior Year Activity Allotment}
Five percent of the student activity fees of each class year, as defined by 
the Constitution of the SGA, shall be set aside into senior year activities 
accounts bound to the graduating class from which the fees were collected. 
This account shall be unavailable for expenditure until the class year to 
which it is bound enters its senior year. The Senior Year Activity Committee, 
by a vote of two-thirds of committee quorum, shall have the power to spend 
from the account once the account's bound year enters its senior year. 

\section{Advertising}
The SGA shall be responsible for informing the Student Body of all campus 
initiatives and for assisting all RSOs in publicizing their events. 

The SGA shall fund flyer requests for all RSOs. Any additional advertising 
requests, as well as any requests by non-RSO Campus Organizations, must be 
approved by the Senate. 

\section{Parliamentary Authority} 
The rules contained in the current edition of Robert's Rules of Order Newly Revised shall serve
as the parliamentary authority for the SGA in which they are not inconsistent with the
Constitution, these bylaws, and any special rules of order the Senate may adopt.


\chapter{Severability}
If a court of competent jurisdiction deems any provision of these bylaws 
unenforceable, that provision will be enforced to the maximum extent 
permissible, and remaining provisions will remain in full force and effect. 

\chapter{Powers Reserved to the Senate} \label{sec:senate_powers}
The Senate shall have the power to carry forth acts when explicitly permitted 
by these bylaws. Any further actions shall require amendment to these bylaws 
to permit these acts. 

\chapter{Powers Reserved to Recognized Student Organizations} \label{sec:rso_powers}
Any powers not explicitly reserved in these bylaws for the Senate and not 
explicitly banned in these bylaws shall be reserved to the Recognized Student 
Organizations.

\end{document}

%%% Local Variables:
%%% mode: latex
%%% TeX-master: t
%%% End:
